\documentclass[12pt]{article}
\usepackage[utf8]{inputenc}
\usepackage[margin=1in]{geometry}
\usepackage{amsmath}\usepackage{amssymb}\usepackage{multicol}\usepackage{enumitem}
\usepackage{pgfplots}\pgfplotsset{compat=1.17}
\setlength{\parindent}{0pt}
\begin{document}

\fbox{\parbox{0.97\linewidth}{\small\textbf{Final answers} (record here --- graded on the answer; show your work):\\[6pt]1.~\underline{\hspace{1.5cm}}\quad 2.~\underline{\hspace{1.5cm}}\quad 3.~\underline{\hspace{1.5cm}}\quad 4.~\underline{\hspace{1.5cm}}\quad 5.~\underline{\hspace{1.5cm}}\quad 6.~\underline{\hspace{1.5cm}}\quad 7.~\underline{\hspace{1.5cm}}\quad 8.~\underline{\hspace{1.5cm}}\quad 9.~\underline{\hspace{1.5cm}}\quad 10.~\underline{\hspace{1.5cm}}}}
\vspace{0.35cm}

\begin{center}{\Large \textbf{State the domain (denominator restrictions)}}\end{center}\vspace{0.2cm}
% State the domain (denominator restrictions)
\textbf{Example (State the domain (denominator restrictions)):}

$f(x)=\dfrac{1}{x+3}$: set $x+3\ne 0$, so $x\ne -3$. Domain: all reals except $-3$.

\vspace{0.2cm}\hrule\vspace{0.35cm}
\textbf{Directions: Set the denominator not equal to zero and solve.}

\begin{multicols}{2}
\begin{enumerate}[label=\textbf{\arabic*.}, start=1, itemsep=1.3cm]
  \item $\displaystyle f(x)=\dfrac{1}{x-4}$
  \item $\displaystyle f(x)=\dfrac{1}{x+7}$
  \item $\displaystyle f(x)=\dfrac{x}{x^{2}-9}$
  \item $\displaystyle f(x)=\dfrac{2}{x^{2}-16}$
\end{enumerate}
\end{multicols}
\vspace{0.25cm}\hrule\vspace{0.35cm}

\begin{center}{\Large \textbf{State the domain (radical restrictions)}}\end{center}\vspace{0.2cm}
% State the domain (radical restrictions)
\textbf{Example (State the domain (radical restrictions)):}

$f(x)=\sqrt{x+6}$: set the radicand $\ge 0$: $x+6\ge 0$, so $x\ge -6$.

\vspace{0.2cm}\hrule\vspace{0.35cm}
\textbf{Directions: Set the radicand $\ge 0$ and solve.}

\begin{multicols}{2}
\begin{enumerate}[label=\textbf{\arabic*.}, start=5, itemsep=1.3cm]
  \item $\displaystyle f(x)=\sqrt{x-1}$
  \item $\displaystyle f(x)=\sqrt{2x-8}$
  \item $\displaystyle f(x)=\sqrt{10-x}$
\end{enumerate}
\end{multicols}
\vspace{0.25cm}\hrule\vspace{0.35cm}

\begin{center}{\Large \textbf{State domain AND range}}\end{center}\vspace{0.2cm}
% State domain AND range
\textbf{Example (State domain AND range):}

$f(x)=x^{2}-5$ has no root or denominator, so the domain is all reals; the lowest output is $-5$, so the range is $y\ge -5$.

\vspace{0.2cm}\hrule\vspace{0.35cm}
\textbf{Directions: State the domain, then find the range by identifying the highest or lowest output.}

\begin{multicols}{2}
\begin{enumerate}[label=\textbf{\arabic*.}, start=8, itemsep=2cm]
  \item $\displaystyle f(x)=x^{2}+7$
  \item \textbf{(challenge)}\ $f(x)=\dfrac{\sqrt{x-1}}{x-4}$ --- give the domain only (two restrictions apply at once)
\end{enumerate}
\end{multicols}

\vspace{0.2cm}\hrule\vspace{0.3cm}
\textbf{Challenge} (same sheet):\\[4pt]
\begin{enumerate}[label=\textbf{\arabic*.}, start=10, itemsep=2.2cm]
  \item \textbf{(challenge)}\ Find the domain of $f(x)=\dfrac{\sqrt{x+3}}{x-1}$, showing both restrictions and how they combine.
\end{enumerate}

\newpage

\begin{center}{\Large \textbf{Answer Key --- fully worked}}\end{center}
\vspace{0.25cm}
\textbf{Procedure:} Denominator restrictions: set the denominator $\ne 0$ and solve for the excluded $x$-value(s).; Radical restrictions: set the radicand $\ge 0$ and solve the resulting inequality.; If a function has both a root and a denominator, apply BOTH restrictions and combine them.; For range, find the highest or lowest output the graph reaches..\\[8pt]
\begin{enumerate}[label=\textbf{\arabic*.},itemsep=7pt]
  \item $x\ne 4$
  \item $x\ne -7$
  \item $x\ne \pm 3$ (since $x^2-9=0$ at $x=3$ and $x=-3$)
  \item $x\ne \pm 4$ (since $x^2-16=0$ at $x=4$ and $x=-4$)
  \item $x\ge 1$
  \item $x\ge 4$ (solve $2x-8\ge 0$)
  \item $x\le 10$ (solve $10-x\ge 0$)
  \item Domain: all real numbers; range: $y\ge 7$
  \item Domain: $x\ge 1$ and $x\ne 4$ (the root requires $x\ge 1$; the denominator excludes $x=4$)
  \item Two restrictions apply. Radical: the radicand must be $\ge 0$, so $x+3\ge 0$, giving $x\ge -3$. Denominator: it cannot be $0$, so $x-1\ne 0$, giving $x\ne 1$. Combine them: start with $x\ge -3$, then remove $x=1$. Domain: $x\ge -3$ and $x\ne 1$ (i.e., $[-3,1)\cup(1,\infty)$).
\end{enumerate}

\end{document}